\documentclass[12pt,a4paper,oneside]{article}
\usepackage[brazil]{babel}
\usepackage[utf8]{inputenc}
\usepackage{olymp}

\setcounter{problem}{4}

\begin{document}

\begin{problem}{Notas}{notas}{2}
 
Os professores de INF110 pretendem analisar o desempenho geral da turma com base na média das notas obtidas na próxima prova prática. Ajude-os fazendo um programa para calcular esta média.

%\Notes
%
%Observações, se tiver.

\InputFile

A entrada contém duas linhas: a primeira linha contém somente um número inteiro $N$, que indica o total de alunos que fizeram a prova; a segunda linha contém $N$ valores reais, cada um deles indicando a nota de um dos alunos. Restrições: todas as notas são valores reais entre 0 e 100 e há no máximo 1000 alunos.


\OutputFile

Seu programa deve gerar apenas uma linha de saída, contendo um valor real de duas casas decimais, representando a média das notas dos alunos.

\Example

\begin{example}
\exmp{
5
92 55 62 84 70
}{
72.60
}%
\end{example}

\begin{example}
\exmp{
3
67.1 82.4 65
}{
71.50
}%
\end{example}


%Comentários sobre o exemplo, se necessário.

\end{problem}

\end{document}
